\documentclass[14pt,a4paper]{extarticle}

% ── Кодировки и язык ──────────────────────────────────────────────────────────
\usepackage[T2A]{fontenc}
\usepackage[utf8]{inputenc}
\usepackage[russian,english]{babel}

% ── Поля страницы (ГОСТ) ──────────────────────────────────────────────────────
\usepackage[left=30mm, right=15mm, top=20mm, bottom=20mm]{geometry}

% ── Межстрочный интервал ──────────────────────────────────────────────────────
\usepackage{setspace}
\onehalfspacing

% ── Математика ────────────────────────────────────────────────────────────────
\usepackage{amsmath, amssymb, amsthm}
\usepackage{mathtools}
\usepackage{bm}

% ── Графика ───────────────────────────────────────────────────────────────────
\usepackage{graphicx}
\usepackage{float}
\usepackage{caption}
\usepackage{subcaption}

% ── Таблицы ───────────────────────────────────────────────────────────────────
\usepackage{booktabs}
\usepackage{tabularx}
\usepackage{array}

% ── Код ───────────────────────────────────────────────────────────────────────
\usepackage{listings}
\usepackage{xcolor}

\definecolor{codebg}{rgb}{0.97,0.97,0.97}
\definecolor{codegreen}{rgb}{0.1,0.5,0.1}
\definecolor{codepurple}{rgb}{0.4,0.0,0.6}
\definecolor{codeblue}{rgb}{0.0,0.2,0.6}
\definecolor{codegray}{rgb}{0.4,0.4,0.4}

\lstdefinestyle{cppstyle}{
    backgroundcolor=\color{codebg},
    commentstyle=\color{codegreen}\itshape,
    keywordstyle=\color{codeblue}\bfseries,
    stringstyle=\color{codepurple},
    numberstyle=\tiny\color{codegray},
    basicstyle=\ttfamily\small,
    breakatwhitespace=false,
    breaklines=true,
    captionpos=b,
    keepspaces=true,
    numbers=left,
    numbersep=6pt,
    showspaces=false,
    showstringspaces=false,
    showtabs=false,
    tabsize=4,
    frame=single,
    rulecolor=\color{black!20},
    language=C++,
    inputencoding=utf8,
    extendedchars=true,
    literate={а}{{\cyra}}1 {б}{{\cyrb}}1 {в}{{\cyrv}}1 {г}{{\cyrg}}1
             {д}{{\cyrd}}1 {е}{{\cyre}}1 {ж}{{\cyrzh}}1 {з}{{\cyrz}}1
             {и}{{\cyri}}1 {й}{{\cyrishrt}}1 {к}{{\cyrk}}1 {л}{{\cyrl}}1
             {м}{{\cyrm}}1 {н}{{\cyrn}}1 {о}{{\cyro}}1 {п}{{\cyrp}}1
             {р}{{\cyrr}}1 {с}{{\cyrs}}1 {т}{{\cyrt}}1 {у}{{\cyru}}1
             {ф}{{\cyrf}}1 {х}{{\cyrh}}1 {ц}{{\cyrc}}1 {ч}{{\cyrch}}1
             {ш}{{\cyrsh}}1 {щ}{{\cyrshch}}1 {ъ}{{\cyrhrdsn}}1 {ы}{{\cyrery}}1
             {ь}{{\cyrsftsn}}1 {э}{{\cyrerev}}1 {ю}{{\cyryu}}1 {я}{{\cyrya}}1
             {А}{{\CYRA}}1 {Б}{{\CYRB}}1 {В}{{\CYRV}}1 {Г}{{\CYRG}}1
             {Д}{{\CYRD}}1 {Е}{{\CYRE}}1 {Ж}{{\CYRZH}}1 {З}{{\CYRZ}}1
             {И}{{\CYRI}}1 {Й}{{\CYRISHRT}}1 {К}{{\CYRK}}1 {Л}{{\CYRL}}1
             {М}{{\CYRM}}1 {Н}{{\CYRN}}1 {О}{{\CYRO}}1 {П}{{\CYRP}}1
             {Р}{{\CYRR}}1 {С}{{\CYRS}}1 {Т}{{\CYRT}}1 {У}{{\CYRU}}1
             {Ф}{{\CYRF}}1 {Х}{{\CYRH}}1 {Ц}{{\CYRC}}1 {Ч}{{\CYRCH}}1
             {Ш}{{\CYRSH}}1 {Щ}{{\CYRSHCH}}1 {Ъ}{{\CYRHRDSN}}1 {Ы}{{\CYRERY}}1
             {Ь}{{\CYRSFTSN}}1 {Э}{{\CYREREV}}1 {Ю}{{\CYRYU}}1 {Я}{{\CYRYA}}1
}

\lstdefinestyle{bashstyle}{
    backgroundcolor=\color{black!90},
    basicstyle=\ttfamily\small\color{white},
    breaklines=true,
    frame=none,
    numbers=none,
    showstringspaces=false,
}

% ── Гиперссылки и цитирование ─────────────────────────────────────────────────
\usepackage[
    colorlinks=true,
    linkcolor=black,
    citecolor=blue,
    urlcolor=blue,
    bookmarks=true
]{hyperref}
\usepackage{cite}

% ── Нумерация секций и содержание ─────────────────────────────────────────────
\usepackage{titlesec}
\usepackage{tocloft}

\titleformat{\section}{\large\bfseries}{\thesection}{1em}{}
\titleformat{\subsection}{\normalsize\bfseries}{\thesubsection}{1em}{}
\titleformat{\subsubsection}{\normalsize\itshape}{\thesubsubsection}{1em}{}

\renewcommand{\cftsecleader}{\cftdotfill{\cftdotsep}}

% ── Переносы и отступы ────────────────────────────────────────────────────────
\usepackage{indentfirst}
\setlength{\parindent}{1.25cm}

% ── Нумерация страниц ─────────────────────────────────────────────────────────
\usepackage{fancyhdr}
\setlength{\headheight}{17pt}
\addtolength{\topmargin}{-3pt}
\pagestyle{fancy}
\fancyhf{}
\fancyhead[R]{\thepage}
\renewcommand{\headrulewidth}{0pt}

\fancypagestyle{plain}{%
    \fancyhf{}
    \fancyhead[R]{\thepage}
    \renewcommand{\headrulewidth}{0pt}
}

% ── Разное ────────────────────────────────────────────────────────────────────
\usepackage{enumitem}
\usepackage{microtype}
\usepackage{url}

% ─────────────────────────────────────────────────────────────────────────────
%  ДОКУМЕНТ
% ─────────────────────────────────────────────────────────────────────────────
\begin{document}

% ── Титульный лист ────────────────────────────────────────────────────────────
\begin{titlepage}
    \begin{center}
        \textbf{Санкт-Петербургский политехнический университет Петра Великого}

        \vspace{2cm}

        \large Курсовая работа

        \vspace{0.5cm}

        по дисциплине:

        \vspace{0.3cm}

        \textbf{<<Объектно-ориентированное программирование>>}

        \vspace{0.8cm}

        на тему:

        \vspace{0.3cm}

        \textbf{<<Система визуализации инерциальной навигации на основе полётного контроллера Pixhawk>>}

        \vfill

        \begin{flushright}
            \begin{tabular}{l l}
                Студент:       & Томилов А.\,С. \\[4pt]
                Преподаватель: & Ананьевский М.\,С. \\[4pt]
                Группа:        & 3331506/30401 \\
            \end{tabular}
        \end{flushright}

        \vspace{2cm}

        Санкт-Петербург

        2026
    \end{center}
\end{titlepage}

% ── Содержание ────────────────────────────────────────────────────────────────
\tableofcontents
\newpage

% ─────────────────────────────────────────────────────────────────────────────
\section*{Введение}
\addcontentsline{toc}{section}{Введение}
% ─────────────────────────────────────────────────────────────────────────────

Инерциальная навигация~--- один из фундаментальных методов определения
положения и ориентации объекта в пространстве без использования внешних
сигналов~\cite{woodman2007}. В основе метода лежат данные инерциального
измерительного блока (IMU~--- \textit{Inertial Measurement Unit}),
включающего трёхосевые акселерометр и гироскоп. Интегрируя угловые
скорости и линейные ускорения во времени, можно восстановить полную
навигационную картину: ориентацию, скорость и положение
объекта~\cite{titterton2004}.

Полётный контроллер Pixhawk является промышленным стандартом в области
автономных летательных аппаратов~\cite{px4_docs}. Он оснащён
высококачественным МЭМС-IMU и поддерживает протокол
MAVLink~\cite{mavlink_proto}~--- открытый телеметрический стандарт,
позволяющий в реальном времени получать сырые данные датчиков на внешний
компьютер~\cite{mavlink_docs}.

Целью данной работы является разработка программной системы, которая
осуществляет чтение сырых IMU-данных (сообщения \texttt{RAW\_IMU}
протокола MAVLink) с полётного контроллера Pixhawk и выполняет их
обработку в реальном времени: вычисление ориентации объекта через
кватернионное интегрирование гироскопа~\cite{kuipers1999}, а также
попытку оценки скорости и позиции посредством двойного интегрирования
ускорения~\cite{groves2013}. Результаты визуализируются в интерактивном
3D-окне с использованием OpenGL и библиотеки GLFW~\cite{glfw_docs}.

\textbf{Задачи работы:}
\begin{enumerate}[leftmargin=2cm]
    \item реализовать модуль приёма данных по протоколу MAVLink через
    последовательный порт;
    \item реализовать алгоритм калибровки и кватернионного интегрирования
    гироскопических данных;
    \item реализовать алгоритм двойного интегрирования акселерометра для
    оценки позиции;
    \item реализовать механизм ZUPT (\textit{Zero Velocity Update}) для
    частичной компенсации дрейфа;
    \item создать OpenGL-визуализацию навигационного состояния в реальном
    времени;
    \item провести анализ точности полученных результатов.
\end{enumerate}

\newpage

% ─────────────────────────────────────────────────────────────────────────────
\section{Чтение сырых данных с Pixhawk}
% ─────────────────────────────────────────────────────────────────────────────

\subsection{Протокол MAVLink и сообщение RAW\_IMU}

Первым этапом работы стало освоение взаимодействия с полётным
контроллером через протокол MAVLink~\cite{mavlink_proto}. MAVLink
(\textit{Micro Air Vehicle Link})~--- это лёгкий бинарный протокол
сериализации для беспилотных систем, разработанный Лоренцем Майером
в 2009~году~\cite{mavlink_docs}. Протокол широко применяется в
экосистеме PX4 и ArduPilot и поддерживается полётным контроллером
Pixhawk~\cite{px4_docs}.

Для получения сырых показаний датчиков используется сообщение
\texttt{RAW\_IMU} (идентификатор~27 в диалекте \texttt{common})~\cite{mavlink_docs}.
Оно содержит девять целочисленных полей в единицах внутренней шкалы
датчика (LSB~--- \textit{Least Significant Bit}):

\begin{table}[H]
    \centering
    \caption{Поля сообщения MAVLink \texttt{RAW\_IMU}}
    \label{tab:rawimu}
    \begin{tabular}{l l l}
        \toprule
        \textbf{Поле} & \textbf{Датчик} & \textbf{Описание} \\
        \midrule
        \texttt{xacc, yacc, zacc}     & Акселерометр & Ускорение по осям X, Y, Z (LSB) \\
        \texttt{xgyro, ygyro, zgyro}  & Гироскоп     & Угловая скорость по осям (LSB) \\
        \texttt{xmag, ymag, zmag}     & Магнитометр  & Магнитное поле по осям (LSB) \\
        \bottomrule
    \end{tabular}
\end{table}

\subsection{Реализация модуля чтения}

Взаимодействие с Pixhawk реализовано через последовательный порт
\texttt{/dev/ttyACM0} со скоростью 115\,200~бод. Инициализация порта
выполнена через POSIX-интерфейс \texttt{termios}~\cite{posix_termios},
позволяющий настраивать все параметры последовательного интерфейса:
скорость, биты данных, паритет, управление потоком.

Ключевые параметры настройки порта: 8 бит данных (CS8), без паритета
(\textasciitilde PARENB), один стоп-бит (\textasciitilde CSTOPB),
аппаратное управление потоком отключено (\textasciitilde CRTSCTS),
таймаут чтения~--- 1~секунда (\texttt{VTIME}~=~10).

После инициализации порта запускается основной цикл, в котором байты
поступают на вход MAVLink-парсера~\cite{mavlink_docs}. Парсер реализован
в виде конечного автомата: функция \texttt{mavlink\_parse\_char()}
возвращает ненулевое значение только при получении полного валидного
пакета, прошедшего проверку контрольной суммы CRC-16~\cite{mavlink_proto}.

\begin{lstlisting}[style=cppstyle, caption={Настройка последовательного порта через termios}]
int setup_serial_port(const char* port, int baudrate) {
    int fd = open(port, O_RDWR | O_NOCTTY);
    struct termios tty;
    memset(&tty, 0, sizeof(tty));
    tcgetattr(fd, &tty);

    cfsetospeed(&tty, baudrate);
    cfsetispeed(&tty, baudrate);

    tty.c_cflag |= (CLOCAL | CREAD);
    tty.c_cflag &= ~CSIZE;
    tty.c_cflag |= CS8;        // 8 data bits
    tty.c_cflag &= ~PARENB;    // no parity
    tty.c_cflag &= ~CSTOPB;    // 1 stop bit
    tty.c_cflag &= ~CRTSCTS;   // no hardware flow control

    tty.c_lflag &= ~(ICANON | ECHO | ECHOE | ISIG);
    tty.c_iflag &= ~(IXON | IXOFF | IXANY);
    tty.c_oflag &= ~OPOST;

    tty.c_cc[VMIN]  = 0;
    tty.c_cc[VTIME] = 10;  // 1 second timeout

    tcsetattr(fd, TCSANOW, &tty);
    return fd;
}
\end{lstlisting}

\begin{lstlisting}[style=cppstyle, caption={Основной цикл чтения и парсинга MAVLink}]
mavlink_message_t msg;
mavlink_status_t  status;
uint8_t buffer[2048];

while (running) {
    int n = read(serial_fd, buffer, sizeof(buffer));
    if (n > 0) {
        for (int i = 0; i < n; i++) {
            if (mavlink_parse_char(MAVLINK_COMM_0,
                                   buffer[i], &msg, &status)) {
                if (msg.msgid == MAVLINK_MSG_ID_RAW_IMU) {
                    mavlink_raw_imu_t imu;
                    mavlink_msg_raw_imu_decode(&msg, &imu);
                    // store imu.xacc, yacc, zacc, xgyro, ...
                }
            }
        }
    }
}
\end{lstlisting}

\subsection{Запись данных в CSV}

Параллельно с обработкой каждый принятый пакет \texttt{RAW\_IMU}
записывается в CSV-файл с меткой времени высокого разрешения,
получаемой через \texttt{gettimeofday()} с точностью до
микросекунды~\cite{posix_termios}. Принудительный сброс буфера на
диск (\texttt{flush()}) выполняется каждые 50~пакетов~--- это
компромисс между надёжностью сохранения данных и накладными расходами
на операции ввода-вывода.

Формат строки файла:

\begin{center}
\texttt{timestamp, AX\_raw, AY\_raw, AZ\_raw, GX\_raw, GY\_raw, GZ\_raw,
MX\_raw, MY\_raw, MZ\_raw}
\end{center}

\subsection{Результаты первичного чтения}

В результате данного этапа была подтверждена успешная работа модуля
чтения: сообщения \texttt{RAW\_IMU} стабильно поступали с частотой
порядка 100~Гц, что соответствует штатной конфигурации IMU на
Pixhawk~\cite{px4_docs}. Сырые данные акселерометра в состоянии
покоя показали устойчивое значение по оси~Z (компонента силы
тяжести), а гироскоп демонстрировал дрейф нуля, подлежащий
компенсации на следующем этапе.

Данные были успешно записаны в CSV-файл для последующего анализа
и подбора калибровочных коэффициентов.

\newpage

% ─────────────────────────────────────────────────────────────────────────────
\section{Архитектура системы визуализации}
% ─────────────────────────────────────────────────────────────────────────────

\subsection{Обзор проекта}

На основе опыта, полученного при реализации модуля чтения, была
разработана полноценная система \textbf{Pixhawk IMU Visualizer}~---
многопоточное приложение на C++, в котором два параллельных потока
выполнения взаимодействуют через потокобезопасный разделяемый буфер.

Первый поток реализует чтение данных: он устанавливает соединение с
Pixhawk через последовательный порт \texttt{/dev/ttyACM0} и непрерывно
парсит входящий поток байтов с помощью MAVLink-парсера~\cite{mavlink_docs}.
При получении полного пакета \texttt{RAW\_IMU} данные помещаются в очередь
и записываются в CSV-файл для последующего анализа.

Второй поток выполняет навигационные вычисления и рендеринг: он извлекает
пакеты из очереди, запускает алгоритм интегрирования и обновляет
трёхмерную сцену на каждом кадре~\cite{glfw_docs}.

\subsection{Структура файлов проекта}

\begin{table}[H]
    \centering
    \caption{Файлы проекта}
    \label{tab:files}
    \begin{tabularx}{\textwidth}{l X}
        \toprule
        \textbf{Файл} & \textbf{Назначение} \\
        \midrule
        \texttt{src/imu\_shared.h} & Общие структуры данных: \texttt{RawIMU}, \texttt{NavState}, \texttt{SharedBuffer}, константы калибровки \\
        \texttt{src/imu\_visualizer.h} & Объявление функции \texttt{run\_visualizer()} \\
        \texttt{src/imu\_visualizer.cpp} & Навигационный алгоритм и OpenGL-рендер \\
        \texttt{src/pixhawk\_imu\_reader.cpp} & MAVLink-ридер, запуск потока визуализации \\
        \texttt{CMakeLists.txt} & Система сборки CMake \\
        \texttt{c\_library\_v2/} & Заголовочная библиотека MAVLink C~\cite{mavlink_docs} \\
        \bottomrule
    \end{tabularx}
\end{table}

\subsection{Межпоточное взаимодействие}

Обмен данными между потоками организован через структуру
\texttt{SharedBuffer}, включающую:

\begin{itemize}[leftmargin=2cm]
    \item \texttt{std::deque<RawIMU> queue}~--- очередь сырых пакетов
    (максимум 500 элементов) с защитой через \texttt{std::mutex};
    \item \texttt{NavState nav}~--- текущее навигационное состояние
    (кватернион, скорость, позиция, углы Эйлера) с отдельным мьютексом;
    \item \texttt{std::deque<array<float,3>> trail}~--- история траектории
    (до 2000 точек) для трёхмерного отображения;
    \item \texttt{std::atomic<bool> running}~--- флаг завершения работы,
    обеспечивающий корректное закрытие обоих потоков.
\end{itemize}

\newpage

% ─────────────────────────────────────────────────────────────────────────────
\section{Зависимости и система сборки}
% ─────────────────────────────────────────────────────────────────────────────

\subsection{Используемые библиотеки}

\begin{table}[H]
    \centering
    \caption{Внешние зависимости проекта}
    \label{tab:deps}
    \begin{tabularx}{\textwidth}{l l X}
        \toprule
        \textbf{Библиотека} & \textbf{Версия} & \textbf{Роль в проекте} \\
        \midrule
        MAVLink \texttt{c\_library\_v2}~\cite{mavlink_docs} & 2.0 & Парсинг телеметрических пакетов, декодирование \texttt{RAW\_IMU} \\
        GLFW3~\cite{glfw_docs} & $\geq$3.3 & Создание окна OpenGL, ввод мыши и клавиатуры \\
        GLEW & $\geq$2.1 & Загрузка расширений OpenGL 3.3 \\
        GLM & $\geq$0.9.9 & Матричная и кватернионная математика в шейдерах \\
        pthreads & POSIX & Межпоточная синхронизация (\texttt{std::thread}) \\
        \bottomrule
    \end{tabularx}
\end{table}

Библиотека MAVLink является заголовочной~--- она полностью состоит из
\texttt{.h}-файлов и не требует отдельной компиляции~\cite{mavlink_proto}.
Остальные зависимости устанавливаются из репозитория пакетов Ubuntu.

\subsection{Конфигурация CMake}

Сборочная система использует CMake версии не ниже 3.10. Скрипт
\texttt{CMakeLists.txt} настроен на стандарт C++14, включает поиск всех
зависимостей через \texttt{find\_package} и \texttt{PkgConfig} и
транслирует проект с флагами оптимизации \texttt{-O2} и предупреждениями
\texttt{-Wall -Wextra}.

\newpage

% ─────────────────────────────────────────────────────────────────────────────
\section{Математика инерциальной навигации}
% ─────────────────────────────────────────────────────────────────────────────

\subsection{Калибровка датчиков}

Сырые показания МЭМС-датчиков содержат постоянное смещение нуля (bias),
обусловленное конструктивными особенностями кристалла~\cite{titterton2004}.
Смещения определялись экспериментально по данным датчика в состоянии
покоя. Перевод в физические единицы осуществляется умножением на
масштабный коэффициент:

\begin{equation}
    a_i = \bigl(\tilde{a}_i - b_{a,i}\bigr) \cdot S_a, \quad
    \omega_i = \bigl(\tilde{\omega}_i - b_{\omega,i}\bigr) \cdot S_\omega,
    \label{eq:calib}
\end{equation}

где $\tilde{a}_i$, $\tilde{\omega}_i$~--- сырые целочисленные отсчёты;
$b_{a,i}$, $b_{\omega,i}$~--- смещения нуля; $S_a = 9{,}81 / 3050$
[м/с$^2$/отсчёт], $S_\omega = 0{,}001$ [рад/с/отсчёт].

Числовые значения смещений, использованных в проекте:

\begin{equation}
    \mathbf{b}_a = (31,\; 24,\; 0)^\top, \quad
    \mathbf{b}_\omega = (-7,\; 1{,}5,\; 3)^\top.
    \label{eq:biases}
\end{equation}

\subsection{Представление ориентации кватернионами}

Ориентация тела в пространстве описывается единичным
кватернионом~\cite{kuipers1999}:

\begin{equation}
    \mathbf{q} = q_w + q_x\,\mathbf{i} + q_y\,\mathbf{j} + q_z\,\mathbf{k},
    \quad \|\mathbf{q}\| = 1.
    \label{eq:quat}
\end{equation}

Кватернионы предпочтительнее матриц вращения и углов Эйлера по ряду
причин~\cite{kuipers1999}: они не страдают от шарнирного замка
(\textit{gimbal lock}), компактны (4 параметра против 9) и численно
стабильны при нормализации.

\subsection{Интегрирование кватерниона}

Кинематическое уравнение для кватерниона связывает его производную с
угловой скоростью $\bm{\omega} = (\omega_x, \omega_y, \omega_z)^\top$ в
системе координат тела~\cite{mahony2008}:

\begin{equation}
    \dot{\mathbf{q}} = \frac{1}{2}\,\mathbf{q} \otimes
    \begin{pmatrix} 0 \\ \bm{\omega} \end{pmatrix},
    \label{eq:qdot}
\end{equation}

где $\otimes$~--- операция умножения кватернионов.

Для численного интегрирования на шаге $\Delta t$ используется точная
формула через ось-угол~\cite{titterton2004}:

\begin{equation}
    \theta = \|\bm{\omega}\| \cdot \Delta t, \qquad
    \Delta\mathbf{q} = \left(\cos\frac{\theta}{2},\;
    \frac{\bm{\omega}}{\|\bm{\omega}\|}\sin\frac{\theta}{2}\right).
    \label{eq:dq}
\end{equation}

Обновление кватерниона и нормализация:

\begin{equation}
    \mathbf{q} \leftarrow \frac{\mathbf{q} \otimes \Delta\mathbf{q}}
    {\|\mathbf{q} \otimes \Delta\mathbf{q}\|}.
    \label{eq:qupdate}
\end{equation}

Нормализация на каждом шаге предотвращает накопление численной ошибки и
гарантирует, что $\mathbf{q}$ остаётся единичным~\cite{mahony2008}.

\subsection{Матрица поворота из кватерниона}

Для перевода векторов из системы координат тела в мировую систему
координат используется матрица вращения $\mathbf{R} \in SO(3)$,
выраженная через компоненты кватерниона~\cite{kuipers1999}:

\begin{equation}
    \mathbf{R} =
    \begin{pmatrix}
        1 - 2(q_y^2 + q_z^2) & 2(q_x q_y - q_w q_z) & 2(q_x q_z + q_w q_y) \\
        2(q_x q_y + q_w q_z) & 1 - 2(q_x^2 + q_z^2) & 2(q_y q_z - q_w q_x) \\
        2(q_x q_z - q_w q_y) & 2(q_y q_z + q_w q_x) & 1 - 2(q_x^2 + q_y^2)
    \end{pmatrix}.
    \label{eq:rot}
\end{equation}

\subsection{Углы Эйлера (конвенция ZYX)}

Из кватерниона вычисляются углы крена $\phi$, тангажа $\theta$ и
рыскания $\psi$ по формулам~\cite{titterton2004}:

\begin{align}
    \phi   &= \operatorname{atan2}\!\bigl(2(q_w q_x + q_y q_z),\;
               1 - 2(q_x^2 + q_y^2)\bigr), \label{eq:roll} \\
    \theta &= \arcsin\!\bigl(2(q_w q_y - q_z q_x)\bigr), \label{eq:pitch} \\
    \psi   &= \operatorname{atan2}\!\bigl(2(q_w q_z + q_x q_y),\;
               1 - 2(q_y^2 + q_z^2)\bigr). \label{eq:yaw}
\end{align}

Функция $\operatorname{atan2}$ обеспечивает однозначный результат в
диапазоне $(-\pi, \pi]$, в отличие от $\arctan$.

\subsection{Двойное интегрирование (скорость и позиция)}

После получения кватерниона ориентации вектор ускорения переводится из
системы координат тела в мировую систему и из него вычитается вектор
гравитационного ускорения~\cite{groves2013}:

\begin{equation}
    \mathbf{a}_\text{world} = \mathbf{R}\,\mathbf{a}_\text{body} -
    \mathbf{g}_\text{world}, \quad
    \mathbf{g}_\text{world} = (0,\; 0,\; -9{,}81)^\top \text{ [м/с}^2\text{]}.
    \label{eq:aworld}
\end{equation}

Последовательным интегрированием методом Эйлера первого порядка
вычисляются скорость и позиция~\cite{woodman2007}:

\begin{align}
    \mathbf{v}(t + \Delta t) &= \mathbf{v}(t) +
    \mathbf{a}_\text{world}(t) \cdot \Delta t, \label{eq:vel} \\
    \mathbf{p}(t + \Delta t) &= \mathbf{p}(t) +
    \mathbf{v}(t) \cdot \Delta t. \label{eq:pos}
\end{align}

\subsection{Метод ZUPT}

Ключевой причиной дрейфа позиции при чистом интегрировании является
накопление ошибок акселерометра~\cite{woodman2007}. Метод ZUPT
(\textit{Zero Velocity Update})~--- классический способ частичной
коррекции INS в моменты покоя платформы~\cite{skog2010}. Если норма
угловой скорости опускается ниже порогового значения, система считает
объект покоящимся и мягко обнуляет скорость:

\begin{equation}
    \text{если } \|\bm{\omega}\| < \omega_\text{th}:
    \quad \mathbf{v} \leftarrow \mathbf{v} \cdot \alpha_\text{damp},
    \label{eq:zupt}
\end{equation}

где $\omega_\text{th} = 0{,}015$ рад/с~--- порог обнаружения покоя,
$\alpha_\text{damp} = 0{,}85$~--- коэффициент затухания. Дополнительно,
при выполнении условия ZUPT и норме мирового ускорения ниже
$0{,}3$ м/с$^2$, ускорение обнуляется полностью.

\newpage

% ─────────────────────────────────────────────────────────────────────────────
\section{Реализация}
% ─────────────────────────────────────────────────────────────────────────────

\subsection{Алгоритм интегрирования}

Ниже приведён фрагмент, реализующий кватернионное интегрирование в
соответствии с формулами~\eqref{eq:dq}--\eqref{eq:qupdate}:

\begin{lstlisting}[style=cppstyle, caption={Кватернионное интегрирование гироскопа}]
double theta = std::sqrt(gx*gx + gy*gy + gz*gz) * dt;
if (theta > 1e-10) {
    double norm_w = std::sqrt(gx*gx + gy*gy + gz*gz);
    double s  = std::sin(theta * 0.5) / norm_w;
    double c  = std::cos(theta * 0.5);
    double dqw = c, dqx = gx*s, dqy = gy*s, dqz = gz*s;

    // q_new = q (x) dq  (quaternion product)
    double qw = n.qw*dqw - n.qx*dqx - n.qy*dqy - n.qz*dqz;
    double qx = n.qw*dqx + n.qx*dqw + n.qy*dqz - n.qz*dqy;
    double qy = n.qw*dqy - n.qx*dqz + n.qy*dqw + n.qz*dqx;
    double qz = n.qw*dqz + n.qx*dqy - n.qy*dqx + n.qz*dqw;

    double norm = std::sqrt(qw*qw + qx*qx + qy*qy + qz*qz);
    n.qw = qw/norm; n.qx = qx/norm;
    n.qy = qy/norm; n.qz = qz/norm;
}
\end{lstlisting}

\subsection{Структуры данных}

Все разделяемые данные инкапсулированы в структурах:

\begin{lstlisting}[style=cppstyle, caption={Ключевые структуры данных (imu\_shared.h)}]
struct NavState {
    double qw=1, qx=0, qy=0, qz=0;  // quaternion
    double wx=0, wy=0, wz=0;          // angular velocity (rad/s)
    double vx=0, vy=0, vz=0;          // velocity (m/s)
    double px=0, py=0, pz=0;          // position (m)
    double roll=0, pitch=0, yaw=0;    // Euler angles (rad)
    double ax=0, ay=0, az=0;          // acceleration (m/s^2)
};
\end{lstlisting}

\subsection{Визуализация}

Интерфейс пользователя разделён на две области экрана~\cite{glfw_docs}.
В верхней части располагается трёхмерная сцена: модель полётного
контроллера, ориентированная в соответствии с текущим кватернионом,
траектория движения в виде синей линии с затуханием и мировые оси
координат. В нижней части~--- HUD с четырьмя временными графиками:
данные акселерометра, гироскопа, углы Эйлера и текущая скорость.

Управление камерой осуществляется через мышь, сброс траектории~---
клавишей \textbf{R}.

\newpage

% ─────────────────────────────────────────────────────────────────────────────
\section{Инструкция по установке и запуску}
% ─────────────────────────────────────────────────────────────────────────────

\subsection{Установка зависимостей}

Все системные зависимости доступны в стандартном репозитории Ubuntu:

\begin{lstlisting}[style=bashstyle, language=bash]
sudo apt update
sudo apt install build-essential cmake git \
    libglfw3-dev libglew-dev libglm-dev libgl1-mesa-dev
\end{lstlisting}

Библиотека MAVLink поставляется в виде заголовочных файлов и
клонируется непосредственно в каталог проекта~\cite{mavlink_docs}:

\begin{lstlisting}[style=bashstyle, language=bash]
cd pixhawk_imu_visualizer
git clone https://github.com/mavlink/c_library_v2.git
\end{lstlisting}

\subsection{Сборка проекта}

\begin{lstlisting}[style=bashstyle, language=bash]
mkdir build && cd build
cmake ..
make -j$(nproc)
\end{lstlisting}

При успешной сборке в каталоге \texttt{build/} появляется исполняемый
файл \texttt{pixhawk\_imu\_reader}.

\subsection{Настройка прав доступа к порту}

Для обращения к последовательному порту без привилегий суперпользователя
необходимо добавить текущего пользователя в группу
\texttt{dialout}~\cite{posix_termios}:

\begin{lstlisting}[style=bashstyle, language=bash]
sudo usermod -a -G dialout $USER
# Re-login required after this command
# Or temporarily (without re-login):
sudo chmod 666 /dev/ttyACM0
\end{lstlisting}

\subsection{Запуск}

Перед запуском необходимо подключить Pixhawk к компьютеру через USB.
Контроллер определяется системой как устройство \texttt{/dev/ttyACM0}~\cite{px4_docs}.

\begin{lstlisting}[style=bashstyle, language=bash]
cd build
./pixhawk_imu_reader
# If permission denied on serial port:
sudo ./pixhawk_imu_reader
\end{lstlisting}

Окно визуализации открывается немедленно при запуске. Данные начинают
поступать по мере инициализации Pixhawk. Для завершения программы
используется сочетание клавиш Ctrl+C в терминале.

\subsection{Управление в окне визуализации}

\begin{table}[H]
    \centering
    \caption{Управление в окне визуализации}
    \label{tab:controls}
    \begin{tabular}{l l}
        \toprule
        \textbf{Действие} & \textbf{Эффект} \\
        \midrule
        ЛКМ + перетаскивание & Вращение камеры вокруг объекта \\
        Колесо мыши & Приближение / отдаление \\
        Клавиша \textbf{R} & Сброс позиции и траектории в $(0,0,0)$ \\
        \textbf{Ctrl+C} & Завершение программы \\
        \bottomrule
    \end{tabular}
\end{table}

\newpage

% ─────────────────────────────────────────────────────────────────────────────
\section{Результаты и анализ}
% ─────────────────────────────────────────────────────────────────────────────

\subsection{Результаты измерения ориентации}

В ходе тестирования была подтверждена корректная работа алгоритма
кватернионного интегрирования. При вращении полётного контроллера по
всем трём осям трёхмерная модель в окне визуализации воспроизводила
движения с высокой точностью. Углы Эйлера, вычисленные из кватерниона
по формулам~\eqref{eq:roll}--\eqref{eq:yaw}, соответствовали физическим
поворотам устройства. Данный результат согласуется с известными
свойствами алгоритма интегрирования через ось-угол~\cite{mahony2008}.

\begin{figure}[H]
    \centering
    \fbox{\parbox{0.85\textwidth}{\centering\vspace{2cm}
        [Скриншот окна визуализации: трёхмерная модель Pixhawk
        при повороте на 45° по оси X]
    \vspace{2cm}}}
    \caption{Визуализация ориентации полётного контроллера}
    \label{fig:orientation}
\end{figure}

\subsection{Результаты оценки позиции}

Алгоритм двойного интегрирования ускорения показал неудовлетворительные
результаты в части оценки реальных координат. Ключевая проблема~---
\textbf{дрейф позиции}: уже через 5--10~секунд движения накопленная
ошибка составляла метры, не соответствующие реальному перемещению. Данная
проблема является фундаментальной для бескоррекционного инерциального
счисления~\cite{woodman2007}.

\begin{figure}[H]
    \centering
    \fbox{\parbox{0.85\textwidth}{\centering\vspace{2cm}
        [График дрейфа позиции: сравнение оценки алгоритма
        с реальным положением за 30 секунд]
    \vspace{2cm}}}
    \caption{Иллюстрация дрейфа позиции при двойном интегрировании}
    \label{fig:drift}
\end{figure}

\subsection{Анализ причин дрейфа}

Дрейф позиции при чистом инерциальном счислении~--- известная и
фундаментальная проблема~\cite{titterton2004,woodman2007}. Рассмотрим
её математически. Пусть истинное ускорение $a(t)$ и ускорение, измеренное
датчиком, $\hat{a}(t)$ связаны:

\begin{equation}
    \hat{a}(t) = a(t) + b + \eta(t),
    \label{eq:noise}
\end{equation}

где $b$~--- постоянное смещение (bias) после калибровки, $\eta(t)$~---
случайный шум. Тогда ошибка позиции после двойного интегрирования:

\begin{equation}
    \delta p(T) = \frac{1}{2}\,b\,T^2 + \int_0^T\!\!(T - t)\,\eta(t)\,dt.
    \label{eq:drift}
\end{equation}

Первое слагаемое растёт квадратично по времени даже при малом остаточном
bias~\cite{groves2013}. Второе слагаемое описывает случайное блуждание.
МЭМС-акселерометры класса Pixhawk имеют типичный шум
$10^{-3}$--$10^{-2}$~м/с$^2$, что за 10~секунд даёт ошибку позиции
порядка единиц метров~\cite{woodman2007}.

\subsection{Методы устранения дрейфа}

Полноценная инерциальная навигация невозможна без внешних корректирующих
сенсоров~\cite{groves2013}. Возможные пути решения проблемы:

\begin{enumerate}[leftmargin=2cm]
    \item \textbf{Фильтр Калмана} с GPS-коррекцией~\cite{groves2013}~---
    классический подход в авиационных навигационных системах. GPS
    предоставляет абсолютную позицию с частотой 1--10~Гц, IMU обеспечивает
    высокую частоту обновления (100--1000~Гц).

    \item \textbf{Фильтр Маджвика}~\cite{madgwick2011}~---
    градиентный алгоритм, сочетающий гироскоп, акселерометр и
    магнитометр. Позволяет существенно улучшить оценку ориентации при
    медленных движениях, однако не решает проблему позиции.

    \item \textbf{Нелинейный дополняющий фильтр на SO(3)}~\cite{mahony2008}~---
    альтернативный подход к оценке ориентации, основанный на теории
    управления на группах Ли.

    \item \textbf{Оптический поток}~\cite{px4_docs}~--- использование
    видеопотока с камеры для оценки скорости относительно поверхности.
    Применяется в системах, где GPS недоступен (помещения).

    \item \textbf{ZUPT с фильтром Калмана}~\cite{skog2010}~---
    расширенное применение метода нулевой скорости в рамках
    оптимального оценщика состояния.
\end{enumerate}

\newpage

% ─────────────────────────────────────────────────────────────────────────────
\section*{Заключение}
\addcontentsline{toc}{section}{Заключение}
% ─────────────────────────────────────────────────────────────────────────────

В ходе выполнения данной работы была разработана система визуализации
инерциальной навигации на основе полётного контроллера Pixhawk.
Реализованы следующие компоненты:

\begin{itemize}[leftmargin=2cm]
    \item модуль приёма и парсинга сообщений \texttt{RAW\_IMU} протокола
    MAVLink~\cite{mavlink_docs} через последовательный интерфейс;
    \item алгоритм калибровки сырых показаний акселерометра и гироскопа;
    \item точное кватернионное интегрирование угловой скорости~\cite{kuipers1999}
    с нормализацией на каждом шаге;
    \item вычисление матрицы вращения и углов Эйлера из кватерниона;
    \item алгоритм двойного интегрирования ускорения для оценки
    позиции~\cite{groves2013};
    \item механизм ZUPT~\cite{skog2010} для частичной компенсации дрейфа скорости;
    \item OpenGL~3.3-визуализация ориентации, траектории и графиков
    датчиков в реальном времени~\cite{glfw_docs}.
\end{itemize}

Алгоритм оценки ориентации показал высокое качество: повороты по всем
трём осям воспроизводились корректно. Алгоритм оценки позиции через
двойное интегрирование ускорения продемонстрировал принципиальную
проблему накопления дрейфа, характерную для МЭМС-датчиков потребительского
класса~\cite{titterton2004,woodman2007}: ошибка позиции нарастает
квадратично по времени и через 5--10 секунд движения достигает метров.

Таким образом, подтверждён известный вывод~\cite{groves2013}: получение
пригодных для практического использования координат возможно лишь при
совместной обработке данных IMU с корректирующими сенсорами в рамках
алгоритма фильтра Калмана или аналогичного оптимального оценщика.
Данная работа создаёт основу для дальнейшего развития системы в этом
направлении.

\newpage

% ─────────────────────────────────────────────────────────────────────────────
\section*{Список использованных источников}
\addcontentsline{toc}{section}{Список использованных источников}
% ─────────────────────────────────────────────────────────────────────────────

\begin{thebibliography}{99}

\bibitem{titterton2004}
    Titterton~D., Weston~J.
    \textit{Strapdown Inertial Navigation Technology}.
    2nd~ed.~--- IET, 2004.~--- 558~p.

\bibitem{woodman2007}
    Woodman~O.\,J.
    \textit{An Introduction to Inertial Navigation}.
    Technical Report UCAM-CL-TR-696.~---
    University of Cambridge, 2007.~--- 37~p.~---
    URL: \url{https://www.cl.cam.ac.uk/techreports/UCAM-CL-TR-696.pdf}
    (дата обр.~01.03.2024).

\bibitem{kuipers1999}
    Kuipers~J.\,B.
    \textit{Quaternions and Rotation Sequences}.~---
    Princeton University Press, 1999.~--- 398~p.

\bibitem{mahony2008}
    Mahony~R., Hamel~T., Pflimlin~J.-M.
    Nonlinear Complementary Filters on the Special Orthogonal Group~//
    \textit{IEEE Transactions on Automatic Control}.~---
    2008.~--- Vol.~53, No.~5.~--- P.~1203--1218.~---
    DOI: \href{https://doi.org/10.1109/TAC.2008.923738}{10.1109/TAC.2008.923738}.

\bibitem{madgwick2011}
    Madgwick~S.\,O.\,H., Harrison~A.\,J.\,L., Vaidyanathan~R.
    Estimation of IMU and MARG orientation using a gradient descent
    algorithm~// \textit{IEEE International Conference on Rehabilitation
    Robotics (ICORR)}.~--- 2011.~--- P.~1--7.~---
    DOI: \href{https://doi.org/10.1109/ICORR.2011.5975346}{10.1109/ICORR.2011.5975346}.

\bibitem{groves2013}
    Groves~P.\,D.
    \textit{Principles of GNSS, Inertial, and Multisensor Integrated
    Navigation Systems}.
    2nd~ed.~--- Artech House, 2013.~--- 776~p.

\bibitem{skog2010}
    Skog~I., Handel~P.
    In-Car Positioning and Navigation Technologies~---
    A Survey~//
    \textit{IEEE Transactions on Intelligent Transportation Systems}.~---
    2009.~--- Vol.~10, No.~1.~--- P.~4--21.~---
    DOI: \href{https://doi.org/10.1109/TITS.2008.2011712}{10.1109/TITS.2008.2011712}.

\bibitem{mavlink_proto}
    Meier~L. et al.
    MAVLink: A Lightweight Message Marshalling Library for the
    DroneCode Ecosystem.~---
    URL: \url{https://github.com/mavlink/mavlink}
    (дата обр.~01.03.2024).

\bibitem{mavlink_docs}
    MAVLink Developer Guide: Common Messages.~---
    URL: \url{https://mavlink.io/en/messages/common.html}
    (дата обр.~01.03.2024).

\bibitem{px4_docs}
    PX4 Autopilot User Guide.~---
    URL: \url{https://docs.px4.io/main/en/}
    (дата обр.~05.03.2024).

\bibitem{glfw_docs}
    GLFW~--- An OpenGL library. Documentation.~---
    URL: \url{https://www.glfw.org/docs/latest/}
    (дата обр.~05.03.2024).

\bibitem{posix_termios}
    The Open Group. The Single UNIX Specification, Version~4.
    \textit{termios.h~--- define values for termios}.~---
    URL: \url{https://pubs.opengroup.org/onlinepubs/9699919799/basedefs/termios.h.html}
    (дата обр.~01.03.2024).

\end{thebibliography}

\end{document}ы